\subsection{埃尔米特几何。}

定义 2. --- 设 A 是一个体，L 是 A 上的一个仿射空间，T 是 L 的平移空间（第 II 章，第 2 \(^{e}\) 版，附录 II）。若在 T 上配备一个非退化埃尔米特型 Φ，则称 L 为 A 上的埃尔米特空间，并称 Φ 为 L 的度量型。

若 J 是恒等映射（这蕴含 A 是交换的），则改称 L 为欧几里得空间。

设 \(a\) 与 \(b\) 是 L 的两个点，令 \(e(a, b) = \Phi(b - a, b - a)\) 。设 \(c\) 是 L 的第三个点。依 § 1 第 5 小节的公式 (17)，要使 \(b - a\) 与 \(c - a\) 正交，必须有 \(e(b, c) = e(a, b) + e(a, c)\) ，而当 J = 1 且 A 的特征不是 2 时，这个条件是充分的（「勾股定理」）。

称 L 的两个线性簇是正交的，如果它们的方向（第 II 章，第 2 \(^{e}\) 版，附录 II，第 3 小节）正交。称 L 的一个线性簇是迷向的（相应地，全迷向的），如果它的方向是迷向的（相应地，全迷向的）。称 T 的一个向量与 L 的一个线性簇正交，如果它与该簇的方向正交。

设 V 是 L 中的一个线性簇，x 是 L 的一个点。L 中使得 y − x 与 V 正交的点 y 所成的集是一个过 x 的线性簇 W；称 W 为过 \(x\) 而与 V 完全正交（或更简单地说，正交）的簇。若 L 是有限维的，则 W 的维数等于 V 的余维数。此外，若 V 是非迷向的，则 V 与 W 的方向互补（§ 4 第 1 小节命题 1 的推论）；于是 W 与 V 交于单独一个点 \(x_1\) ；在 V 中取定一个原点，立即可以看出，对固定的 V，映射 \(x \to x_1\) 是一个幂等仿射线性映射；称之为 L 到 V 上的正交投影；与之相伴的线性映射（第 II 章，第 2 \(^{e}\) 版，附录 II，第 4 小节）是 T 到 V 的方向上的正交投影算子（第 3 小节）。

定义 3. --- 设 L 是体 A 上的一个埃尔米特空间，T 是 L 的平移空间。称 L 的位移（相应地，相似变换）为 L 到 L 上的任一仿射双射 u，使得 u 在 T 中与之相伴的线性映射 v（第 II 章，第 2 \(^{e}\) 版，附录 II，第 4 小节）是酉的（相应地，是一个相似变换）。

平移群是仿射群的正规子群；因而它是相似变换群的正规子群，也是位移群的正规子群。对每个 \(a \in L\) ，以 \(G_a\) 表示使 \(a\) 不动的相似变换（相应地，位移）所成的群；若取 \(a\) 为原点而把 L 等同于 T，则 \(G_a\) 是 T 的相似变换群（相应地，酉群）。每个相似变换（相应地，位移）\(u\) 都可以以唯一的方式写成 \(u = u_1 t_1\) 的形式，其中 \(u_1 \in G_a\) 且 \(t_1 \in T\) ，也可以写成 \(u = t_2 u_2\) 的形式，其中 \(u_2 \in G_a\) 且 \(t_2 \in T\) ；此外，有 \(u_2 = u_1\) 及 \(t_2 = u_1 t_1 u_1^{-1}\) （第 II 章，第 2 \(^\text{e}\) 版，附录 II，第 4 小节）。

设 \(u\) 是 L 中的一个相似变换，\(\varrho\) 是 T 中与之相伴的相似变换。\(\varrho\) 的乘子也称为 \(u\) 的乘子（第 5 小节）。若以 \(\alpha(u)\) 记这个乘子，则映射 \(u \to \alpha(u)\) 是从 L 的相似变换群到 A 的中心的可逆元素乘法群中的同态；其核是位移群，因而后者是相似变换群的正规子群。当 A 交换且 L 有限维时，在行列式 det \(u\) （依定义等于 det \(\varrho\) ）与 \(\alpha(u)\) 之间，成立与第 5 小节中相同的关系。诸位移 \(u\) ，凡使 det \(u = 1\) 者，构成位移群的一个正规子群；若 A 是特征 ≠ 2 的域且 J 是恒等映射，则这个子群的指标为 2。

命题 6. --- 设 L 是 A 上一个有限维埃尔米特空间，其度量型的指数为 0。则 L 的每个乘子为 \(\mu \neq 1\) 的相似变换 u 都有且只有一个不动点。

实际上，设 \(a\) 是 L 的一个点。存在 L 的相似变换 \(\nu\) 使 \(a\) 不动，又存在 L 的平移 \(t\) ，使得 \(u = tv\) 。说 \(b\) 是 \(u\) 的不动点，等于说 \(\nu(b) - b = t\) 。为证明这个方程有且只有一个解 \(b\) ，取 \(a\) 为原点而把 L 与其平移空间 T 等同。于是只需证明 T 的自同态 \(\nu - 1\) 可逆，换言之，证明关系 \(\nu(x) - x = 0\) ( \(x \in \mathrm{T}\) ) 蕴含 \(x = 0\) 。而今，若 \(\nu(x) - x = 0\) ，则有 \(\Phi(x, x) = \Phi(\nu(x), \nu(x)) = \mu \Phi(x, x)\) ，故得 \(\Phi(x, x) = 0\) ，因为 \(\mu \neq 1\) ；这蕴含 \(x = 0\) ，因为 \(\Phi\) 的指数为 0。证毕。

设 A 是特征 \(\neq 2\) 的体。L 的每个位移 \(u\) ，若满足 \(u^2 = 1\) ，便至少有一个不动点，例如两个对应点的中点 \(\frac{1}{2}(x + u(x))\) ；取这个点为原点，可见与 \(u\) 相伴的 T 的酉自同构是一个对称（第 3 小节）。设 V 是 L 中的一个非迷向线性簇；称位移 \(u\) 为关于 V 的对称，如果在 V 中取定一个原点后，\(u\) 就等同于 T 关于 V 的对称。这等价于说 \(u(x)\) 如下得到：以 \(x_1\) 表示 \(x\) 在 V 上的正交投影，我们有 \(u(x) - x = 2(x_1 - x)\) 。

习题. -- 1) 设 A 是域。给定 A 上一个 n 阶埃尔米特矩阵 R，称 R 的 r 阶主子式为从 R 中删去 n − r 行以及具有相同指标的 n − r 列所得到的诸子式。

\begin{enumerate}
\def\labelenumi{\alph{enumi})}
\setcounter{enumi}{1}
\tightlist
\item
  由 a) 推出：若 R 的秩为 r，则存在一个置换 \(\sigma \in S_{n}\) ，使得：当把同一置换 \(\sigma\) 施行于 R 的行与列，以 S 记这样得到的矩阵，并以 \(\Delta_{k}\) 记从 S 中删去指标 \textgreater{} k 的行与列所得的 S 的 k 阶主子式时，有下列两条性质：\(1^{\circ} \Delta_{r} \neq 0\) ；\(2^{\circ}\) 不存在指标 k \textless{} r 使 \(\Delta_{k} = \Delta_{k+1} = 0\) 。
\end{enumerate}

\begin{enumerate}
\def\labelenumi{\arabic{enumi})}
\setcounter{enumi}{1}
\tightlist
\item
  假设 A 是域，且 E 的维数为有限数 n。设 \(\Phi\) 是 E 上的埃尔米特半双线性型，满足 § 4 第 2 小节的条件 (T)，\(R = (\alpha_{ij})\) 是 \(\Phi\) 关于 E 的基 ( \(e_i\) ) 的矩阵。
\end{enumerate}

\begin{enumerate}
\def\labelenumi{\alph{enumi})}
\item
  若 \(\Phi\) 的秩为 \(r\) ，且从 \(R\) 中删去指标 \(>r\) 的行与列所得的主子式（习题 1）非零，证明：存在 E 的新基 \((f_i)\) ，使得 \(e_i = f_i\) （对 \(1 \leqslant i \leqslant r\) ），且 \(\Phi\) 关于 \((f_i)\) 的矩阵是在 \(R\) 中把诸 \(\alpha_{ij}\) （ \(i > r\) 或 \(j > r\) ）换成 0 所得的矩阵（考虑与 E 正交的子空间 E \(^\circ\) ）。
\item
  由 a) 推出：若 \(\Phi\) 的秩为 \(n\) ，且 \(\Delta_{n-1}\) （ \(\alpha_{nn}\) 在行列式 \(\Delta = \det R\) 中的余子式）非零，则存在 E 的新基 \((f_i)\) ，使得 \(f_i = e_i\) （对 \(1 \leqslant i \leqslant n-1\) ），并且有
\end{enumerate}

\[
\Phi (x, y) = \Phi \left(\sum_ {i = 1} ^ {n} \xi_ {i} f _ {i}, \sum_ {i = 1} ^ {n} \eta_ {i} f _ {i}\right) = \sum_ {i = 1} ^ {n - 1} \sum_ {j = 1} ^ {n - 1} \alpha_ {i j} \xi_ {i} \bar {\eta} _ {j} + \frac {\Delta}{\Delta_ {n - 1}} \xi_ {n} \bar {\eta} _ {n}
\]

（考虑这样的埃尔米特型：其关于 \((e_i)\) 的矩阵是把 \(\alpha_{nn}\) 换成 \(\alpha_{nn} - \frac{\Delta}{\Delta_{n-1}}\) 后在 \(R\) 中所得的矩阵）。

\begin{enumerate}
\def\labelenumi{\alph{enumi})}
\setcounter{enumi}{2}
\tightlist
\item
  假设 \(\Phi\) 的秩为 \(n\) ，\(\Delta_{n-1} = 0\) ，但主子式 \(\Delta_{n-2}\) （由在 \(R\) 中删去指标 \(n-1\) 与 \(n\) 所在的行与列而得的 \(R\) 的主子式）非零。证明：存在 E 的新基 \((f_i)\) ，使得 \(f_i = e_i\) （对 \(1 \leqslant i \leqslant n-2\) ），并且有
\end{enumerate}

\[
\Phi (x, y) = (\sum_ {i = 1} ^ {n} \xi_ {i} f _ {i}, \sum_ {i = 1} ^ {n} \eta_ {i} f _ {i}) = \sum_ {i = 1} ^ {n - 2} \sum_ {j = 1} ^ {n - 2} \alpha_ {i j} \xi_ {i} \bar {\eta} _ {j} + \xi_ {n - 1} \bar {\eta} _ {n} + \xi_ {n} \bar {\eta} _ {n - 1}.
\]

（若 H 是由 \(e_{1},\ldots,e_{n-1}\) 生成的超平面且它是迷向的，注意与 H 正交的直线不含于由 \(e_{1},\ldots,e_{n-2}\) 生成的子空间，并利用 § 4 第 2 小节命题 2）。

\begin{enumerate}
\def\labelenumi{\arabic{enumi})}
\setcounter{enumi}{2}
\item
  设 A 是有限域，E 是 A 上有限维向量空间，Φ 是 E 上的非退化埃尔米特半双线性型，关于 A 的一个异于恒等映射的自同构 J。证明 E 具有一个关于 Φ 的标准正交基（参见第 V 章 § 11 第 5 小节定理 3 的推论）。
\item
  设 A 是特征 \(\neq 2\) 的有限域，E 是 A 上维数为 \(n\) 的有限维向量空间。
\end{enumerate}

\begin{enumerate}
\def\labelenumi{\alph{enumi})}
\item
  证明：对 E 上每个非退化对称双线性型 \(\Phi\) ，存在 E 的正交基 \((e_i)\) ，使得 \(\Phi(e_i, e_i) = 1\) （对 \(1 \leqslant i \leqslant n - 1\) ），且 \(\Phi(e_n, e_n) = \Delta\) （\(\Phi\) 关于 \((e_i)\) 的判别式）。（注意若 \(\alpha\beta \neq 0\) ，则方程 \(\alpha\xi^2 + \beta\eta^2 = \gamma\) 在 A 中总有解 \((\xi, \eta)\) （若 \(\gamma \neq 0\) ）（第 V 章 § 11 习题 4）。）
\item
  要使 E 上两个非退化对称双线性型等价，必须且只需它们的判别式（关于 E 的同一基取的）之比是 A 中的一个平方。由此推出：若 n 是奇数，则对 E 上每个非退化对称双线性型 \(\Phi\) ，存在一个正交基，使 \(\Phi\) 关于它的矩阵具有 \(\lambda I_{n} (\lambda \in \mathbf{A})\) 的形式；此时 \(\Phi\) 的指数为 \((n - 1)/2\) 。
\item
  若 \(n=2m\) 是偶数，证明：E 上非退化对称双线性型 \(\Phi\) 的指数为 \(m\) （若 \((-1)^{m}\Delta\) 是 A 中的平方），否则为 \(m-1\) 。
\end{enumerate}

\begin{enumerate}
\def\labelenumi{\arabic{enumi})}
\setcounter{enumi}{4}
\tightlist
\item
  设 A 是特征 \(\neq 2\) 的域。设 I 是系数在 A 中的多项式，含 \(n(n+1)/2\) 个未定元 \(X_{ij}\) ( \(1 \leqslant i \leqslant j \leqslant n\) )；对 A 的交换扩域 A′ 上的每个对称矩阵 \(R = (\alpha_{ij})\) ，以 I(R) 表示把 \(\alpha_{ij}\) 代入 I 中以代替未定元 \(X_{ij}\) ( \(i \leqslant j\) ) 所得的 A′ 的元素。
\end{enumerate}

假设 I 具有如下性质：对于矩阵 \(U = (u_{ij})\) （其中 \(u_{ij} = \mathrm{X}_{ij}\) （对 \(i \leqslant j\) ），\(u_{ij} = \mathrm{X}_{ji}\) （对 \(i > j\) ））以及方阵 \(P = (\mathrm{Y}_{ij})\) （ \(n\) 阶；其中诸 \(\mathrm{Y}_{ij}\) 是另外 \(n^2\) 个未定元），有

\[
\mathrm{I} (\mathbf {t} P U P) = (\det P) ^ {\hbar} \mathrm{I} (U)
\]

其中 h 是一个整数 \textgreater{} 0。证明 h 是偶数，且 I(U) = γ(det U)k，其中 h = 2k 且 γ ∈ A。（利用定理 1 证明：对 A 的代数闭包 Ω 上的每个对称矩阵 R，有 (I(R))² = λ(det R)h，其中 λ ∈ Ω；并利用 det U 关于诸 Xij 的多项式不是一个平方这一事实，这可由考察此多项式中含某个 Xii 的各项得出。）

*6) 设 A 是完备非离散赋值域，交换且特征 \(\neq 2\) （《一般拓扑》，第 IX 章，§ 3，第 2 小节），\(\Phi\) 是 A 上维数为 \(n\) 的有限维向量空间 E 上的非退化埃尔米特型，\(R = (\alpha_{ij})\) 是 \(\Phi\) 关于基 \((e_i)\) 的矩阵。证明存在 \(\varepsilon > 0\) ，使得对每个埃尔米特矩阵 \(R' = (\alpha_{ij}')\) ，只要它满足条件 \(|\alpha_{ij}' - \alpha_{ij}| \leqslant \varepsilon\) （对每对 \((i, j)\) ），型 \(\Phi'\) （以 \(R'\) 为关于基 \((e_i)\) 的矩阵）便与 \(\Phi\) 等价。（归结为 \(R\) 是对角矩阵的情形；利用习题 2 b)，并依赖下述引理：存在一个数 \(a > 0\) ，使当 \(|\eta| \leqslant a\) 时，A 中存在元素 \(\xi\) 满足 \(\xi^2 = 1 - \eta\) 。为证明这个引理，利用 \((1 - x)^{1/2})\cdot *\)

¶ 7) 设 A 是特征 ≠ 2 的不可序域（第 VI 章 § 2 习题 8），E 是 A 上维数 n \textgreater{} 0 的有限维向量空间，Q 是 E 上的非退化二次型，\((e_i)\) 是关于 Q 的正交基，使得 \(\mathrm{Q}\left(\sum_{i=1}^{n}\xi_{i}e_{i}\right)=\sum_{i=1}^{n}\alpha_{i}\xi_{i}^{2}\) 。对 \(1\leqslant r\leqslant n\) ，令 \(\mathrm{Q}_{r}(\xi_{1},\ldots,\xi_{r})=\sum_{i=1}^{r}\alpha_{i}\xi_{i}^{2}\) ，并以 \(M_{r}\) 表示 \(Q_{r}\) 的值当诸 \(\xi_{i}\left(1\leqslant i\leqslant r\right)\) 取遍 A 时所成的集。

\begin{enumerate}
\def\labelenumi{\alph{enumi})}
\item
  证明：若对某个指标 \(r\) 有 \(\mathbf{M}_r = \mathbf{M}_{r+1}\) ，则由此推出 \(\mathbf{M}_r = \mathbf{A}\) （注意 A 的每个元素都是一些平方之和（第 VI 章 § 2 习题 7））。
\item
  假设 A* 的由 A 中元素的平方所成的子群 S 在 A* 中有有限指标 s。由 a) 推出：若 n \textgreater{} s，则 E 上每个非退化二次型的指数都 \textgreater{} 0 （注意每个集 M \(_{r}\) 都是 0 与模 S 的若干类之并）。*应用于 A 是 p 进域 Q \(_{p}\) 的情形（《一般拓扑》，第 III 章，§ 5，习题 35）。*
\end{enumerate}

\begin{enumerate}
\def\labelenumi{\arabic{enumi})}
\setcounter{enumi}{7}
\item
  设 A 是特征 \(\neq 2\) 的域，E 是 A 上维数为 \(n\) 的有限维向量空间，Q 是 E 上指数为 0 的非退化二次型。设 A′ 是 A 的一个有限奇次数代数扩张，E′ 是把 E 的标量域扩张到 A′ 所得的 A′ 上的向量空间。证明 Q 到 E′ 上的扩张 Q′（§ 3 第 4 小节命题 3）仍具有指数 0。（归结为 A′ = A{[}X{]}/(f) 的情形，其中 f 是 A 上次数为奇数 m 的不可约多项式。设 ( \(e_i\) ) 是 E 关于 Q 的正交基，且 \(\rho_i = Q(e_i)\) ；证明：在 A{[}X{]} 中，形如 \(\sum_{i} \rho_i(g_i(X))^2 = f(X)h(X)\) 的关系（其中诸 \(g_i\) 是次数 \(\leqslant m - 1\) 且不全为零的多项式）是不可能的；为此注意 h 必有奇次数，并考察 h 的一个次数为奇数的不可约因子）。
\item
  设 A 是一个体，E 是 A 上具有可数基 \((e_n)_{n\geqslant 1}\) 的向量空间，\(\Phi\) 是 E 上的非退化埃尔米特半双线性型，满足条件 (T)（ \(\S 4\) ，第 2 小节）。
\end{enumerate}

\begin{enumerate}
\def\labelenumi{\alph{enumi})}
\item
  证明：若定理 1 的条件 (C) 不同时满足，则 E 中存在关于 \(\Phi\) 的正交基（按 § 5 习题 4 的方式推理）。
\item
  再假设 A 是交换的，且存在整数 s，使得在 A 上每个维数 \textgreater s 的有限维向量空间上，每个非退化埃尔米特半双线性型的指数都 \textgreater0 （参见习题 7）。证明此时 E 中存在关于 Φ 的标准正交基。（按 a) 的方式推理，注意对 A 中每个形如 α = λ + λ̄ 的元素，以及维数 \textgreater s 的有限维空间 F 上的每个非退化埃尔米特型 Ψ，存在 z ∈ F 使 Ψ(z, z) = α（参见 § 4 第 2 小节命题 4）。）
\end{enumerate}

¶ 10) a) 设 A 是主理想环，其中只有一个极大理想 Aπ，且 2 不被 π 整除（第 VII 章 § 1 习题 4）。设 E 是 A 上维数为 n 的自由模。证明 E 上每个对称双线性型 Φ 都具有正交基。（设 r 是使 πr 整除所有元素 Φ(x, y) 的最大指数；证明存在 a ∈ E 使 Φ(a, a) = απr，其中 α 在 A 中可逆；由此推出 E 是 F = Aa 与 F 的正交子模 F⁰ 的直和。）

\begin{enumerate}
\def\labelenumi{\alph{enumi})}
\setcounter{enumi}{1}
\item
  举一个例子（对 n = 2），其中 \(\Phi\) 非退化，且存在 E 的一个非迷向子模 F，秩为 1，在 E 中有补，但 \(F^{0}\) 不是 F 的补。
\item
  设 \((e_i)\) 是关于 \(\Phi\) 的正交基，且 \(\alpha_i = \Phi(e_i, e_i)\) 。证明：诸理想 \(\mathbf{A}\alpha_i\) 在不计次序的意义下与所考虑的正交基无关（参见 § 5 定理 1）。
\end{enumerate}

称这些理想为型 \(\Phi\) 的不变因子。举两个具有相同不变因子而不等价的型的例子（取判别式之比不是平方的两个型）。

\begin{enumerate}
\def\labelenumi{\alph{enumi})}
\setcounter{enumi}{3}
\item
  设 F 是 E 的子模，\(\Phi_{\mathrm{F}}\) 是 \(\Phi\) 在 \(\mathbf{F} \times \mathbf{F}\) 上的限制，\(\mathbf{A}\alpha_{i}\) ( \(1 \leqslant i \leqslant r\) ) 是 \(\Phi\) 的非零不变因子，排列得使 \(\alpha_{i}\) 整除 \(\alpha_{i+1}\) ，\(\mathrm{A}\beta_{i}\) ( \(1 \leqslant i \leqslant s\) ) 是 \(\Phi_{\mathrm{F}}\) 的非零不变因子，排列得使 \(\beta_{i}\) 整除 \(\beta_{i+1}\) 。证明有 \(s \leqslant r\) ，且 \(\beta_{i}\) 是 \(\alpha_{i}\) 的倍式（对 \(1 \leqslant i \leqslant s\) ）（方法同 § 5 习题 1 a)）。
\item
  假设 \(\Phi\) 非退化；设 F、G 是 E 的两个非迷向子模，使得 \(F^{0}\) （相应地 \(G^{0}\) ）是 F（相应地 G）的补。假设 \(\Phi\) 在 F 与 G 上的限制等价；证明此时存在 E 的自同构 \(u\) ，保持 \(\Phi\) 不变，且使 \(u(F) = G\) 。（利用 \(a\) )，可归结为 \(F = Aa\) 、\(G = Ab\) 、\(\Phi(a, a) = \Phi(b, b)\) 的情形。设 \((c_{j})\) 是 \(G^{0}\) 的一个基，且设 \(b'\) 、\(c_{j}'\) ( \(1 \leqslant j \leqslant n - 1\) ) 分别是 \(b\) 与诸 \(c_{j}\) 在 \(F^{0}\) 中的分量；证明存在标量 \(\mu_{j}\) ( \(1 \leqslant j \leqslant n - 1\) )，使元素 \(d_{j} = c_{j}' + \mu_{j}b'\) 对每对指标满足关系 \(\Phi(d_{j}, d_{k}) = \Phi(c_{j}, c_{k})\) ；为此注意，对每个 \(\lambda \in A\) ，元素 \(1 \pm \lambda\) 之一在 A 中可逆。）
\end{enumerate}

\begin{enumerate}
\def\labelenumi{\arabic{enumi})}
\setcounter{enumi}{10}
\item
  设 A 是特征 0 的主理想环，其中只有一个主理想 \(\pi\) ，使得 2 被 \(\pi\) 整除。若 \((e_1, e_2)\) 是 \(E = A^2\) 的标准基，\(\Phi\) 是 E 上由 \(\Phi(\xi_1 e_1 + \xi_2 e_2, \eta_1 e_1 + \eta_2 e_2) = \xi_1 \eta_2 + \xi_2 \eta_1\) 定义的对称双线性型，证明不存在 E 的关于 \(\Phi\) 的正交基。
\item
  设 A 是有限域 \(\mathbf{F}_{q^{2}}\) ，J 是 A 的对合自同构 \(\xi \to \xi^{q}\) ，其固定域是 \(\mathbf{F}_{q}\) 。若 E 是 A 上维数为 \(n\) 的向量空间，\(\Phi\) 是 E 上（关于 J 的）非退化埃尔米特半双线性型，证明酉群 \(\mathbf{U}(\Phi)\) 的阶等于
\end{enumerate}

\[
(q ^ {n} - (- 1) ^ {n}) q ^ {n - 1} (q _ {-} ^ {n - 1} - (- 1) ^ {n - 1}) q ^ {n - 2} \dots (q ^ {2} - 1) q (q + 1)
\]

（用与 § 5 习题 10 类似的方法，并借助习题 3）。

\begin{enumerate}
\def\labelenumi{\arabic{enumi})}
\setcounter{enumi}{12}
\tightlist
\item
  设 A 是有限域 \(F_{q}\) （q 不是 2 的倍数），E 是 A 上维数为 n 的向量空间，Q 是 E 上的非退化二次型。证明：
\end{enumerate}

\begin{enumerate}
\def\labelenumi{\alph{enumi})}
\tightlist
\item
  若 n 是奇数，群 SO(Q) 的阶为
\end{enumerate}

\[
(q ^ {n - 1} - 1) q ^ {n - 2} (q ^ {n - 3} - 1) q ^ {n - 4} \dots (q ^ {2} - 1) q.
\]

\begin{enumerate}
\def\labelenumi{\alph{enumi})}
\setcounter{enumi}{1}
\tightlist
\item
  若 n = 2m 是偶数，群 SO(Q) 的阶等于
\end{enumerate}

\[
(q ^ {2 m - 1} - \varepsilon q ^ {m - 1}) (q ^ {2 m - 2} - 1) q ^ {2 m - 3} \dots (q ^ {2} - 1) q
\]

其中 \(\varepsilon = 1\) 若 \((-1)^{m}\Delta\) 是 A 中的平方，否则 \(\varepsilon = -1\) ，\(\Delta\) 表示 Q 关于 E 的任意基的判别式。（方法类似于习题 12，利用 § 6 的习题 3 及第 V 章 § 11 的习题 5。）

\begin{enumerate}
\def\labelenumi{\arabic{enumi})}
\setcounter{enumi}{13}
\tightlist
\item
  假设 A 是域，E 是 A 上维数 \(n \geqslant 2\) 的有限维向量空间，\(\Phi\) 是 E 上的非退化埃尔米特半双线性型，满足条件 (T)（ \(\S 4\) ，第 2 小节）。证明：E 的自同态 \(w\) ，凡与所有自同构 \(u\) （属于特殊酉群 \(\mathbf{SU}(\Phi)\) 者）交换的，只有位似，除非同时有 \(n = 2\) 、J = 1 且 A 的特征为 \(\neq 2\) 。（若 \(n \geqslant 3\) ，写出 \(w\) 与诸对合 \(u \in \mathbf{SU}(\Phi)\) 交换，并利用 \(\S 4\) 的习题 3；若 \(n = 2\) 且 J \(\neq 1\) ，写出 \(w\) 与特殊酉群 \(\mathbf{SU}(\Phi)\) 中那些关于 E 的一个正交基的矩阵形如 \(\begin{pmatrix} \lambda & 0 \\ 0 & \lambda^{-1} \end{pmatrix}\) 的元素交换。
\end{enumerate}

¶ 15) 设 A 是特征 ≠ 2 的域，E 是 A 上维数 n ≥ 1 的向量空间，Q 是 E 上的非退化二次型。对每个自同构 u ∈ O(Q)，令 w = u − 1，以 r 记 w 的秩，并令 W = ̄w(0)。

\begin{enumerate}
\def\labelenumi{\alph{enumi})}
\item
  证明 \(w(E)\) 是 \(W^{0}\) （与 \(W\) 正交的子空间）。
\item
  证明：若 \(n=2\) 、\(r=2\) ，则 \(u\) 是关于 E 中两条直线的两个对称之积。（证明：若 \(w(x)\) 对每个非迷向向量 \(x\in\mathrm{E}\) 都是迷向的，则 \(w(x)\) 对每个 \(x\in\mathrm{E}\) 都是迷向的；将 A 至少有 5 个元素的情形与 \(\mathrm{A}=\mathbf{F}_{3}\) 的情形分别讨论。）
\item
  设 n 与 r 是任意的。证明：若 \(w(\mathrm{E})\) 不是全迷向的，则 u 是 r 个关于 E 的超平面的对称之积，且不能表为个数更少的对称之积。（归结为 W 全迷向的情形，并对 n 与 r 作归纳。若 \(W \neq \{0\}\) ，证明存在向量 \(a \in W^{0}\) 使 \(w(a)\) 非迷向，用反证法论证，并利用如下事实：一个平面上除至多一条外所有直线都迷向时它必全迷向；然后取关于与 \(w(a)\) 正交的超平面的对称 s，并考察自同构 su。若 \(W = \{0\}\) ，取 \(a \in E\) 使 \(w(a)\) 非迷向，s 的意义同上，再考察自同构 su，并利用 b)。）
\item
  假设 \(w(E)\) 是全迷向的。若 \(s\) 是关于非迷向超平面 H 的一个对称，证明在 \(su\) 下不变的向量所成的子空间是 H ∩ W，故其维数为 \(n-r-1\) ，并由此推出 \(su\) 不能是少于 \(r+1\) 个关于超平面的对称之积。然后由 \(c\) ) 推出 \(u\) 是 \(r+2\) 个关于超平面的对称之积，但不能是个数更少的对称之积。
\item
  由 c) 与 d) 推出：每个正交自同构都是至多 n 个关于超平面的对称之积。
\item
  证明：若 n 是奇数（相应地，偶数），则对每个行列式为 1（相应地，−1）的自同构 \(u \in \mathbf{0}(\mathrm{Q})\) ，存在在 u 下不变的向量 \(x \neq 0\) （利用 e)）。
\end{enumerate}

\begin{enumerate}
\def\labelenumi{\arabic{enumi})}
\setcounter{enumi}{15}
\tightlist
\item
  在与习题 15 相同的假设下，证明：若 \(n \geqslant 3\) ，群 SO(Q) 由关于 E 的维数为 \(n - 2\) 的非迷向子空间的诸对称生成（按第 4 小节命题 5 的方式推理）。
\end{enumerate}

¶ 17) 假设与习题 15 的相同。

\begin{enumerate}
\def\labelenumi{\alph{enumi})}
\item
  证明：对 \(n \geqslant 2\) ，\(\Omega(Q)\) （正交群 \(\mathbf{O}(Q)\) 的换位子群）由元素 \((st)^{2}\) 生成，其中 \(s\) 与 \(t\) 取遍关于超平面的诸对称（利用第 4 小节命题 5，并注意对任意群 \(\Gamma\) ，由 \(\Gamma\) 的元素的平方生成的子群包含 \(\Gamma\) 的换位子群）。
\item
  证明：若 \(n \geqslant 3\) ，\(\mathbf{SO}(\mathbf{Q})\) 的换位子群由 \(\mathbf{SO}(\mathbf{Q})\) 的元素的平方生成（利用习题 16）；由此推出这个群等同于 \(\Omega(\mathbf{Q})\) ，且商群 \(\mathbf{SO}(\mathbf{Q})/\Omega(\mathbf{Q})\) 是一个交换群，其所有元素的阶都是 2。
\item
  称平面 P ⊂ E 为双曲的，如果它非迷向且包含迷向直线（必然有两条）。称自同构 \(u \in \mathbf{0}(Q)\) 为双曲的，如果存在双曲平面 P 使 \(u(x) = x\) （对一切 \(x \in \mathrm{P}^0\) ）；这时称 \(u\) 为与 P 相伴的双曲变换。证明：若 Q 的指数 \(\geqslant 1\) ，则每个 \(u \in \mathbf{0}(Q)\) 都是一些双曲变换之积（利用第 4 小节命题 5 及 § 4 习题 4 a)）。由此推出：若 P 是双曲平面，则每个 \(u \in \mathbf{0}(Q)\) 都可写成 \(u = tv\) ，其中 \(t\) 是与 P 相伴的双曲变换，且 \(v \in \Omega(Q)\) 。
\end{enumerate}

¶ 18) 设 A 是交换域，E 是 A 上维数为 n 的向量空间，Φ 是 E 上的非退化埃尔米特半双线性型，满足条件 (T)（§ 4 第 2 小节）。设 V 是 E 的一个向量子空间，H \(_{v}\) 是酉群 U(Φ) 中由使 u(V) = V 的酉自同构 u 所成的子群。

\begin{enumerate}
\def\labelenumi{\alph{enumi})}
\item
  证明：当 V 不是维数为 \(n/2\) 的全迷向子空间时，\(\mathrm{H}_{\mathrm{v}}\) 在映射 \(u \to \det u\) 下的像是 \(\mathbf{A}^*\) 中由诸 \(\rho \in \mathbf{A}\) （满足 \(\rho \bar{\rho} = 1\) ）组成的子群。
\item
  若 \(n\) 是偶数且 \(V\) 是维数为 \(n/2\) 的全迷向子空间，证明 \(H_{v}\) 在映射 \(u \rightarrow \det u\) 下的像是 A* 中由形如 \(\bar{\lambda}/\lambda\) 的元素所成的子群（利用 § 4 第 2 小节命题 2）。
\item
  设 V、W 是 E 的两个向量子空间，使得 \(\Phi\) 在 V 与 W 上的限制等价。证明在下列情形中存在 \(u \in \mathbf{S}\mathbf{U}(\Phi)\) 使 \(u(\mathrm{V}) = \mathrm{W}\) ：
\end{enumerate}

\(^{1}\) J 异于恒等映射（利用第 V 章 § 11 第 5 小节定理 3）。

\(2^{\circ} \text{ J} = 1, \text{ A 的特征为 } \neq 2, \text{ V 与 W 不是维数为 } n/2.\)

\begin{enumerate}
\def\labelenumi{\alph{enumi})}
\setcounter{enumi}{3}
\tightlist
\item
  假设 J = 1，A 的特征 ≠ 2，n = 2m 是偶数，且 Φ 是指数为 m 的非退化对称双线性型。设 V、W 是 E 中两个维数为 m 的全迷向子空间；证明：若 dim(V ∩ W) = q，则对每个使 u(V) = W 的正交自同构 u，有 det u = \((-1)^{m-q}\) （利用 b) 及 § 4 第 2 小节命题 2）。由此推出：维数为 m 的全迷向子空间所成的集是群 SU(Φ) 的两个非传递类 N \(_{1}\) 、N \(_{2}\) 之并；若 V 与 W 属于同一类（相应地，属于不同的类），则 V ∩ W 的维数与 m 有相同的奇偶性（相应地，与 m 的奇偶性不同）。要使（关于 Φ 的）相似变换 u 是正的，必须且只需 \(u(N_{1}) = N_{1}\) （利用 § 4 习题 4 c)）。
\end{enumerate}

\begin{enumerate}
\def\labelenumi{\arabic{enumi})}
\setcounter{enumi}{18}
\tightlist
\item
  设 A 是一个体，E 是 A 上维数 \textgreater{} 零的有限维向量空间，\(\Phi\) 是 E 上的非退化且非交错的 \(\varepsilon\) -埃尔米特半双线性型。设 u 是 E 的一个自同构，使得
\end{enumerate}

\[
\Phi (u (x), u (x)) = \alpha \Phi (x, x)
\]

对一切 \(x \in \mathbf{E}\) 成立，其中 \(\alpha \in \mathbf{A}\) 。证明 \(u\) 是乘子为 \(\alpha\) 的相似变换，除非下列条件同时成立：A 交换且特征为 2，并且 J 是恒等映射（利用 § 1 习题 8）。

\begin{enumerate}
\def\labelenumi{\arabic{enumi})}
\setcounter{enumi}{19}
\item
  设 A 是一个体，L 是 A 上有限维埃尔米特空间；假设 L 的度量型 \(\Phi\) 满足条件 (T)（ \(\S\) 4，第 2 小节）。证明：若 \(\Phi\) 的指数 \(>0\) ，则 L 的乘子 \(\neq 1\) 且没有不动点的相似变换可能存在（利用第 6 小节命题 6 的推理，及 \(\S\) 4 第 2 小节命题 2）。
\item
  设 A 是特征 ≠ 2 的域，L 是 A 上有限维欧几里得空间，Φ 是 L 的度量型。
\end{enumerate}

\begin{enumerate}
\def\labelenumi{\alph{enumi})}
\tightlist
\item
  证明：L 到自身的每个满足
\end{enumerate}

\[
\Phi (u (x) - u (y), u (x) - u (y)) = \Phi (x - y, x - y)
\]

（对一切 \(x, y\) ∈ \(L\) ）的双射 u 都是一个位移（利用 § 1 习题 7）。

\begin{enumerate}
\def\labelenumi{\alph{enumi})}
\setcounter{enumi}{1}
\tightlist
\item
  证明：位移群由关于仿射空间 L 的非迷向超平面的诸对称生成（利用第 4 小节命题 5，归结为证明每个非迷向平移是两个这种对称之积）。
\end{enumerate}

\begin{enumerate}
\def\labelenumi{\arabic{enumi})}
\setcounter{enumi}{21}
\item
在埃尔米特空间 L 中，称两个线性簇是垂直的，如果它们的方向是弱正交子空间（§ 3 习题 11）。设 L 是有限维的；设 \(V_{1}\) 、\(V_{2}\) 是两个线性簇，\(W_{1}\) 、\(W_{2}\) 是它们各自的方向。假设 \(p = \dim(W_{1} + W_{2}) < n\) ；证明：若 \(W_{1} + W_{2}\) 不是迷向的，则至少存在一个维数为 n - p 的线性簇 U，同时垂直于 \(V_{1}\) 与 \(V_{2}\) ，且与簇 \(V_{1}\) 、\(V_{2}\) 各恰交于一点；此外，若 \(q = \dim(W_{1} \cap W_{2})\) ，则具有前述性质的诸线性簇 U 的并是一个维数为 \(n - p + q\) 的线性簇。
\item
设 A 是特征 \(\neq 2\) 的域，E 是 A 上维数 \(n + 1 \geqslant 2\) 的有限维向量空间，Q 是 E 上的二次型，\(\Phi\) 是与 Q 相伴的对称双线性型。称 \(x \in E\) （满足 \(Q(x) = 0\) ）所成的集 C 为以 0 为顶点、以 \(Q(x) = 0\) 为方程的迷向锥。若它不化为 0，则 \(C - \{0\}\) 在射影空间 \(\mathbf{P}(\mathrm{E})\) 中在典范映射 \(\pi\) （从 \(\mathrm{E}-\{0\}\) 到 \(\mathbf{P}(\mathrm{E})\) 上；第 II 章，\(2^{\mathrm{e}}\) 版，附录 III）下的像 S 称为射影二次曲面（相应地，若 \(n=2\) 则为射影二次曲线），其齐次方程为 \(\mathrm{Q}(x)=0\) 。若 Q 退化，则称 S 是退化的。称射影线性簇 \(\mathrm{V}_{1}, \mathrm{V}_{2}\) （ \(\mathbf{P}(\mathrm{E})\) 中的）关于 S 共轭，如果 \(\overline{\pi}(\mathrm{V}_{1})\) 与 \(\overline{\pi}(\mathrm{V}_{2})\) （关于 \(\Phi\) ）正交。极簇 \(\mathrm{V}^{0}\) （射影线性簇 \(\mathrm{V}\subset\mathbf{P}(\mathrm{E})\) 关于 S 的极簇）是这样的簇：\(\overline{\pi}(\mathrm{V}^{0})\cup\{0\}\) 是（关于 \(\Phi\) ）与 \(\overline{\pi}(\mathrm{V})\cup\{0\}\) 完全正交的子空间；若 V 是超平面且 S 非退化，则 \(\mathrm{V}^{0}\) 化为单独一点，称为 V 的极点。称射影线性簇 V 与 S 相切，如果 \(\overline{\pi}(\mathrm{V})\cup\{0\}\) 是（关于 \(\Phi\) 的）迷向子空间。
\end{enumerate}

我们在以下总假设 S 非空且非退化。

\begin{enumerate}
\def\labelenumi{\alph{enumi})}
\item
  证明：S 与射影线性簇 V 的交或为空，或为 V 中的一个二次曲面；要使这个二次曲面退化，必须且只需 V 与 S 相切。
\item
  证明：S 在点 \(z \in S\) 处的切超平面是过 z 且与 S 相切的诸直线之并。
\end{enumerate}

\[
z ^ {\prime}
\]

z 关于 a 与 b，即直线 D 上满足 \(\begin{bmatrix} a & b \\ z' & z \end{bmatrix} = -1\) 的点（第 II 章，第 2 版，附录 III，习题 4）；证明 \(z'\) 属于 z 关于 S 的极超平面，且这些点中有 n 个构成 P(E) 中的一个射影自由族并属于 S（参见 § 4 习题 4 a)）。

\begin{enumerate}
\def\labelenumi{\alph{enumi})}
\setcounter{enumi}{3}
\item
  假设 n = 3 且 \(\Phi\) 有极大指数 v = 2。S 所含直线的集这时是两个集 \(N_{1}\) 、\(N_{2}\) 之并，使得 \(N_{1}\) 的每条直线与 \(N_{2}\) 的每条直线相交，而 \(N_{1}\) （相应地 \(N_{2}\) ）的两条不同直线不相交（习题 18 d)）。设 D、\(D'\) 是属于 \(N_{1}\) 的两条不同直线；对每个 \(z \in D\) ，存在恰好一条过 z 的直线 \(\Delta \in N_{2}\) ；若 \(u(z)\) 是 \(\Delta\) 与 \(D'\) 的交点，证明 u 是从 D 到 \(D'\) 上的射影线性映射。
\item
  仍设 n = 3，设 D、D′、D″ 是 P(E) 中三条两两不相交的直线。证明与 D、D′、D″ 都相交的直线的并是一个非退化二次曲面。
\end{enumerate}

\begin{enumerate}
\def\labelenumi{\arabic{enumi})}
\setcounter{enumi}{23}
\tightlist
\item
  假设与记号同习题 23，其中二次曲面 S 假设为非空且非退化。
\end{enumerate}

\begin{enumerate}
\def\labelenumi{\alph{enumi})}
\item
  证明：射影群 PGL(E) 中由把 S 映到自身的诸射影线性双射所成的子群 \(\Gamma\) 是关于 Q 的相似变换群的典范像。（利用上面的习题 23 c)、§ 4 习题 2 a) 及 § 1 习题 8。）
\item
  设 \(a\) 是 \(\mathbf{P}(\mathrm{E})\) 中不属于 S 的一个点，\(\Phi_{1}\) 是 \(\Phi\) 在 E 中与 \(\overline{\pi}(a)\) 正交的超平面上的限制。证明 \(\Gamma\) 中使 \(a\) 不变的子群同构于正交群 \(\mathbf{U}(\Phi_1)\) 关于其中心的商群。
\item
  设 \(b\) 是 S 的一个点，F 是 E 中与 \(\overline{\pi}(b)\) 正交的（迷向）超平面，M 是 \(\overline{\pi}(b)\) 关于 F 的一个（非迷向）补，\(\Phi_{2}\) 是 \(\Phi\) 在 M 上的限制。证明 \(\Gamma\) 中使 \(b\) 不变的子群同构于维数为 \(n-1\) 的欧几里得空间 L 的相似变换群，后者的度量型是（ \(\S\) 1 第 7 小节意义下的） \(\Phi_{2}\) 的逆型。（注意若一个关于 \(\Phi\) 的相似变换把直线 \(\overline{\pi}(b)\) 映到自身，则它把 F 映到自身，且由它在 F 上的限制完全确定。）
\end{enumerate}

\begin{enumerate}
\def\labelenumi{\arabic{enumi})}
\setcounter{enumi}{24}
\tightlist
\item
  设 A 是特征 \(\neq 2\) 的域，L 是 A 上维数 \(n \geqslant 2\) 的有限维仿射空间。我们把 L 等同于射影超平面 \(\mathrm{H}_{0}\) （“无穷远超平面”，属于射影空间 \(\mathbf{P}(\mathrm{E})\) ，后者维数为 \(n\) ；第 II 章，\(2^{\mathrm{e}}\) 版，附录 III，第 4 小节）的补。称非空集 S ⊂ L 为仿射二次曲面（若 \(n = 2\) ，则称仿射二次曲线），如果 S 是 L 与 \(\mathbf{P}(\mathrm{E})\) 中一个射影二次曲面（相应地，二次曲线）的交（习题 23）。
\end{enumerate}

\begin{enumerate}
\def\labelenumi{\alph{enumi})}
\item
  证明：若存在非退化射影二次曲面 \(\overline{\mathrm{S}}\subset\mathbf{P}(\mathrm{E})\) 使 \(S=L\cap\overline{S}\) ，则这个二次曲面是具有这些性质的唯一一个，除非 \(n=2\) 、\(A=F_{3}\) 且 S 化为 2 个元素（注意除这个例外情形外，对每个点 \(z\in H_{0}\) （不属于 \(\overline{S}\) ），存在过 \(z\) 且与 S 交于两个不同点的直线）。这时称 S 为非退化仿射二次曲面。称 L 中所含的两个仿射线性簇 \(V_{1},V_{2}\) 关于 S 共轭，如果射影线性簇 \(\overline{V}_{1},\overline{V}_{2}\) （满足 \(V_{i}=L\cap\overline{V}_{i}\) ( \(i=1,2\) ) 者）关于 \(\overline{S}\) 共轭；类似地定义仿射线性簇关于 S 的极簇（当它不含于 \(H_{0}\) 时）或极点，以及与 S 相切的仿射线性簇。
\item
  假设 S 非退化；证明可在 L 中选取一个原点 a，使得按这种方式把 L 等同于一个向量
\end{enumerate}

空间后，存在 L 的一个基 \((e_i)\) ，使得 S 是满足两种形式之一的方程的 \(x = \sum_{i=1}^{\infty} \xi_i e_i\) 所成的集

\[
\begin{array}{c} \alpha_ {1} \xi_ {1} ^ {2} + \dots + \alpha_ {n} \xi_ {n} ^ {2} = 1 \\ \alpha_ {1} \xi_ {1} ^ {2} + \dots + \alpha_ {n - 1} \xi_ {n - 1} ^ {2} + \xi_ {n} = 0. \end{array}
\]

在第一种情形，点 \(a\) 是完全确定的，且是关于 \(\bar{S}\) 的无穷远超平面 \(H_{0}\) 的极点（称为 S 的中心）。（按 \(H_{0}\) 与 \(\bar{S}\) 相切或不相切区分两种情形；利用 § 6 第 1 小节的定理 1 及 § 4 第 2 小节的命题 2。）

\begin{enumerate}
\def\labelenumi{\arabic{enumi})}
\setcounter{enumi}{25}
\item
  设 A 是代数闭的特征 \(\neq 2\) 的交换域，E 是 A 上有限维向量空间，Q 是 E 上的非退化二次型。设 \(u \in \mathbf{0}(Q)\) ；以 § 4 习题 12 的记号，我们有 G(p, p) = \{0\} ，但 \(p(X) = X - 1\) 与 \(p(X) = X + 1\) 除外。设 M 是含于 \(G(p, p)\) 且在 \(u\) 下稳定的非迷向子空间所成之集中的一个极小元，并设 \(p^h\) 是 \(u\) 在 M 上的限制的极小多项式。证明：若 \(h\) 是奇数，则 M 是 \(E_u\) 的一个不可分解子模；若 \(h\) 是偶数，则 M 是 \(E_u\) 的两个同构的不可分解子模的直和（为看出当 \(h = 2k\) 是偶数时 M 不可能不可分解，证明此时 \(N = p^k(u)(M)\) 将是全迷向的；若 \((e_i)_{1 \leqslant i \leqslant 2k}\) 是 M 的满足 \(u(e_i) = \varepsilon e_i + e_{i+1}\) （对 \(i \leqslant 2k - 1\) ）、\(u(e_{2k}) = \varepsilon e_{2k}\) （其中 \(\varepsilon = \pm 1\) ）的基，证明 \(e_k\) 不能与 \(e_{k+1}\) 正交，并由此推出关系 \(Q(u(e_k)) = Q(e_k)\) 导致矛盾）。
\item
  设 A 是特征 2 的域，E 是 A 上维数为有限数 n 的向量空间，Q 是 E 上的二次型，Φ 是相伴的双线性型，它是交错的，故有偶秩 2m（§ 5 第 1 小节定理 1 的推论 1）。
\end{enumerate}

\begin{enumerate}
\def\labelenumi{\alph{enumi})}
\item
  证明：若 \(E^{0}\) 是 E 的关于 \(\Phi\) 与 E 正交的子空间（维数为 n - 2m），则 \(Q(\lambda x + \mu y) = \lambda^{2}Q(x) + \mu^{2}Q(y)\) 对 \(E^{0}\) 中一切 x、y 成立；换言之，\(Q_{0}\) （Q 在 \(E^{0}\) 上的限制）是从 \(E^{0}\) （视为 A 上的向量空间）到 A（视为子域 \(A^{2}\) 上的向量空间）中的、关于同构 \(\xi \to \xi^{2}\) （从 A 到 \(A^{2}\) 上的同构）的半线性映射。设 q 是 Q 的核 \(E^{0} \cap \overline{Q}(0)\) 的（在 A 上的）维数，并设 \(E_{1}\) 是 \(E^{0} \cap \overline{Q}(0)\) 关于 \(E^{0}\) 的一个补；我们有 \(n - 2m - q \leqslant [A : A^{2}]\) 。
\item
  由 a) 推出：存在 E 的一个基 \((e_i)_{1 \leqslant i \leqslant n}\) ，其前 \(2m\) 个向量构成 \(\mathbf{E}_2\) （\(\mathbf{E}^0\) 在
\end{enumerate}

E 中的一个补）的基，随后的 \(n - 2m - q\) 个向量构成 \(\mathbf{E}_1\) 的基，使得对 \(x = \sum_{i=1}^{n} \xi_i e_i\) 有

\[
\mathrm{Q} (x) = \sum_ {i = 1} ^ {m} \left(\alpha_ {i} \xi_ {i} ^ {2} + \xi_ {i} \xi_ {m + i} + \beta_ {i} \xi_ {m + i} ^ {2}\right) + \sum_ {i = 2 m + 1} ^ {n - q} \gamma_ {i} \xi_ {i} ^ {2}
\]

其中诸 \(\gamma_{i}\) ( \(2m + 1 \leqslant i \leqslant n - q\) ) 是 A 中关于 \(A^{2}\) 线性无关的元素。

\begin{enumerate}
\def\labelenumi{\alph{enumi})}
\setcounter{enumi}{2}
\item
  我们把 Q 的指数定义为使得 \(V \cap E^{0} = \{0\}\) 的全奇异子空间 V 的最大维数。证明：若 v 是 Q 的指数，则可取具有 b) 中所述性质的 E 的基 \((e_{i})\) ，使得 \(\alpha_{i} = \beta_{i} = 0\) （对 \(1 \leqslant i \leqslant v\) ），且 Q 在 \(E_{2}\) 的由 \(e_{v+1}, \ldots, e_{m}, e_{m+v+1}, \ldots, e_{2m}\) 生成的子空间上的限制是指数为 0 的非退化二次型。
\item
  假设 \(q=0\) ；以 \(\mathbf{O}(\mathrm{Q})\) 记使 Q 不变的 E 的自同构所成的群。若 \(u\in\mathbf{O}(\mathrm{Q})\) ，证明 \(u(x)=x\) （对一切 \(x\in\mathrm{E}^{0}\) ）。对每个 \(x\in\mathrm{E}_{2}\) ，令 \(u(x)=u_{0}(x)+u_{2}(x)\) ，其中 \(u_{0}(x)\in\mathrm{E}^{0}\) 且 \(u_{2}(x)\in\mathrm{E}_{2}\) ；证明 \(u_{2}\) 属于辛群 \(\mathbf{Sp}(\Phi_{2})\) （其中 \(\Phi_{2}\) 是 \(\Phi\) 在 \(\mathrm{E}_{2}\) 上的限制），且有 \(\mathrm{Q}(u_{2}(x))+\mathrm{Q}(x)\in\mathrm{Q}(\mathrm{E}^{0})\) 。反之，对每个自同构 \(u_{2}\in\mathbf{Sp}(\Phi_{2})\) ，若它满足 \(\mathrm{Q}(u_{2}(x))+\mathrm{Q}(x)\in\mathrm{Q}(\mathrm{E}^{0})\) （对一切 \(x\in\mathrm{E}_{2}\) ），证明存在唯一的线性映射 \(u_{0}\) ，它从 \(\mathrm{E}_{2}\) 映到 E \(^{0}\) ，使得等于 \(u_{0} + u_{2}\) （在 E \(_{2}\) 上）、等于恒等映射（在 E \(^{0}\) 上）的那个线性映射属于 O(Q)。
\item
  假设 A 是完全域（A² = A）且 q = 0。由 b) 推出：E 的每个维数 ≥ 3 的向量子空间都至少含一个使 Q(x) = 0 的向量 x。若 n 是奇数，则必有 m = v 且 n = 2m + 1，于是存在 E 的基 (ei)，关于它有
\end{enumerate}

\[
\mathrm{Q} \left(\sum_ {i = 1} ^ {n} \xi_ {i} e _ {i}\right) = \xi_ {1} \xi_ {m + 1} + \dots + \xi_ {m} \xi_ {2 m} + \xi_ {2 m + 1} ^ {2},
\]

且（用 \(d\) ) 的记号）\(\mathbf{O}(\mathrm{Q})\) 同构于 \(\mathbf{Sp}(\Phi_{2})\) ；此时所有满足 \(q=0\) 的二次型都等价。若 \(n\) 是偶数，则必有 \(n=2m\) 、\(\nu=m\) 或 \(\nu=m-1\) ，且存在 E 的基 \((e_{i})\) ，关于它有

\[
\mathrm{Q} \left(\sum_ {i = 1} ^ {n} \xi_ {i} e _ {i}\right) = \xi_ {1} \xi_ {m + 1} + \dots + \xi_ {m - 1} \xi_ {2 m - 1} + \xi_ {m} \xi_ {2 m} + \lambda \left(\xi_ {m} ^ {2} + \xi_ {2 m} ^ {2}\right)\tag{1}
\]

其中 \(\lambda\in A\) 。设 \(A_{1}\) 是在 A 上添加多项式 \(\lambda X^{2}+X+\lambda\) 的根所得的域，证明这个域与使 Q 可写成形式 (1) 的基 \((e_{i})\) 无关，且要使两个（满足 \(q=0\) 的）二次型等价，必须且只需它们按此方式对应的 A 的二次扩张相同（利用维特定理）。A 是特征 2 的有限域的情形。

¶ 28) 设 A 是特征 2 的域，异于 \(\mathbf{F}_{2}\) ，E 是 A 上维数为 \(n = 2m\) 的向量空间，Q 是 E 上的非退化二次型。

\begin{enumerate}
\def\labelenumi{\alph{enumi})}
\item
  证明：正交群 \(\mathbf{O}(\mathrm{Q})\) 由诸对称生成（这里它们无非是属于 \(\mathbf{O}(\mathrm{Q})\) 的平延（ \(\S 4\) ，习题 6））（按 \(\S 5\) 习题 11 的方式推理）。由此推出 \(\mathbf{O}(\mathrm{Q})\) 的换位子群由 \(\mathbf{O}(\mathrm{Q})\) 的元素的平方生成（参见习题 17）。
\item
  假设 Q 有极大指数；设 V、W 是 E 的两个维数为 m 的全奇异子空间（§ 4 第 1 小节）。设 u 是对称 \(x \to x + \frac{\Phi(x, a)}{\mathrm{Q}(a)} a\) （§ 4 习题 6）；设 k 是 V ∩ W 的维数。证明 V ∩ u(W) 的维数当 a 与 V ∩ W 正交时为 \(k + 1\) ，在相反的情形为 \(k - 1\) （在第一种情形，注意 \(a = x + y\) ，其中 \(x \in V\) 、\(y \in W\) ，并证明 \(u(y) = x\) ；在第二种情形，注意 u 不能使任何不与 a 正交的奇异向量不变）。
\item
  再假设 Q 的指数是任意的。证明：\(\mathbf{SO}(\mathrm{Q})\) （\(\mathbf{O}(\mathrm{Q})\) 中由 E 的偶数个对称之积的诸自同构所成的子群）是 \(\mathbf{O}(\mathrm{Q})\) 的指标为 2 的正规子群。（证明奇数个对称之积不能是恒等映射，这可通过考察 Q 到
\end{enumerate}

把 E 的标量域扩张到其代数闭包所得的向量空间 E′ 上的扩张来进行；然后利用 b)。）（参见 § 9 习题 9。）

\begin{enumerate}
\def\labelenumi{\alph{enumi})}
\setcounter{enumi}{3}
\item
  若 \(V_{1}\) 、\(\bar{V}_{2}\) 是 E 的两个全奇异子空间，维数相同且 \textless m，证明存在自同构 \(u \in \mathbf{SO}(Q)\) 使 \(u(V_{1}) = V_{2}\) 。反之，若 \(V_{1}\) 与 \(V_{2}\) 是两个维数为 m 的全奇异子空间，则要存在自同构 \(u \in \mathbf{SO}(Q)\) 使 \(u(V_{1}) = V_{2}\) ，必须且只需 \(V_{1} \cap V_{2}\) 的维数与 m 有相同的奇偶性（按习题 18 的方式论证，利用 b)）。
\item
  称平面 P ⊂ E 为双曲的，如果它非迷向且包含奇异直线（必然有两条）。称变换 \(u \in \mathbf{0}(Q)\) 为双曲的，如果存在双曲平面 P 使 \(u(x) = x\) （对一切 \(x \in \mathrm{P}^0\) ）；这时称 \(u\) 为与 P 相伴的双曲变换。证明：若 Q 的指数 \textgreater0，则每个 \(u \in \mathbf{0}(Q)\) 都是一些双曲变换之积（利用 \(a\) )）。由此推出：若 P 是双曲平面，则每个变换 \(u \in \mathbf{0}(Q)\) 都可写成 \(u = sv\) ，其中 \(s\) 是与 P 相伴的双曲变换，且 \(v\) 属于 \(\mathbf{0}(Q)\) 的换位子群。
\end{enumerate}

\begin{enumerate}
\def\labelenumi{\arabic{enumi})}
\setcounter{enumi}{28}
\tightlist
\item
  在习题 28 的假设下，再假设 Q 有极大指数 m；设 \((e_{i})\) 是 E 的由奇异向量组成的（关于与 Q 相伴的交错型 \(\Phi\) 的）辛基（§ 4 第 2 小节命题 2），使 \(\Phi\) 关于这个基的矩阵是 § 5 习题 14 中记作 R 的矩阵。用该习题的记号，证明：要使辛矩阵 \((^{t}D + S)^{-1}(D + S)\) 是某个自同构 \(u \in \mathbf{O}(\mathbb{Q})\) 的矩阵，必须且只需 S 是交错的（写出每个向量 \(u(e_{i})\) 都是奇异的，并注意有 \((^{t}D + S).u(e_{i}) = (D + S).e_{i}\) ）。
\end{enumerate}

\section{埃尔米特型与有序域}

在本节中始终，K 表示一个极大有序域（从而是交换的且特征为零；参见第 VI 章 § 2），并且我们假设处于下列三种情形之一：

1°) A = K，J 是恒等映射；

2°) A 是域 \(\mathrm{K}(i)\) ，它是在 K 上添加 −1 的一个平方根 i 得到的，且对每个 \(\lambda \in A\) ，\(\overline{\lambda}\) 是 \(\lambda\) 的共轭元素（第 II 章 § 7 第 7 小节）。

3°) A 是 K 上对应于偶 \((-1,-1)\) 的四元数代数（或按我们以后的简称，K 上的四元数体），且对每个 \(\lambda\in A\) ，\(\overline{\lambda}\) 是 \(\lambda\) 的共轭元素（参见第 II 章 § 7 第 8 小节及第 VIII 章 § 11 第 2 小节）。

若 \(\Phi\) 是 E 上的埃尔米特型，则在所有情形下，\(\Phi(x, x) \in K\) （对一切 \(x \in E\) ），因为 \(\Phi(x, x) = \overline{\Phi(x, x)}\) 。

\subsection{正埃尔米特型。}

定义 1. --- 称 E 上的埃尔米特型 \(\Phi\) 是正的（相应地，负的），如果 \(\Phi(x, x) \geqslant 0\) （相应地，\(\Phi(x, x) \leqslant 0\) ）对一切 \(x \in \mathbf{E}\) 成立。

当 A = K 时，也称 \(Q(x) = \frac{1}{2} \Phi(x, x)\) （与 \(\Phi\) 相伴的二次型）是正的（相应地，负的）。

设 E 在 A 上有限维，且设 \((e_{i})\) ( \(i = 1, \ldots, n\) ) 是 E 的正交基（ \(\S 6, no. 1\) ，定理 1）。要使 E 上的埃尔米特型 \(\Phi\) 是正的，必须且只需 \(\Phi(e_{i}, e_{i}) \geqslant 0\) （对 \(i = 1, \ldots, n\) ）。设 \(\Phi\) 是 E 上的一个非退化正埃尔米特型；由于 K 的每个正元素都是平方，因而具有形式 \(\rho\bar{\rho} (\rho \in A)\) ，故 E 中存在关于 \(\Phi\) 的标准正交基（ \(\S 6, no. 1\) ，定理 1 的推论 1）。

命题 1. --- 假设 E 有限维，且 A = K 或 A = K(i)。若 \(\Phi\) 是 E 上的非退化正埃尔米特型，则 \(\Phi\) 到 \(\bigotimes^{p}\mathbf{E}\) 与 \(\bigwedge^{p}\mathbf{E}\) ( \(p > 0\) ) 上的扩张以及 \(\Phi\) 的逆型都是非退化正埃尔米特型。

这由 E 的标准正交基的存在性及 § 6 第 1 小节命题 2 立即得出。

若把各处的“非退化正”都换成“正”，命题 1 仍成立。

命题 2. --- 设 Φ 是 E 上的一个正埃尔米特型。对 E 中的 x、y，有

\[
\Phi (x, y) \overline {{{{\Phi (x , y)}}}} \leqslant \Phi (x, x) \Phi (y, y).\tag{1}
\]

当向量 \(x\) 与 \(y\) 成比例时，不等式显然是直接的。因此设 \(x\) 与 \(y\) 线性无关。设 A′ 是 A 的（交换）子域 K( \(\Phi(x, y)\) )，F 是 A′ 上由 \(x\) 与 \(y\) 生成的向量平面，\(\Phi_{\mathbb{F}}\) 是 \(\Phi\) 在 F 上的限制；它取值于 A′ 中。由命题 1，\(\Phi_{\mathbb{F}}\) 关于基 \((x, y)\) 的判别式 \(\geqslant 0\) 。而这个判别式就是 \(\Phi(x, x)\Phi(y, y) - \Phi(x, y)\overline{\Phi(x, y)}\) 。证毕。

推论. --- E 的迷向向量所成的集是子空间 \(E^{0}\) （E 关于 \(\Phi\) 的正交子空间）。要使 \(\Phi\) 非退化，必须且只需 \(\Phi(x, x) > 0\) （对一切 \(x \neq 0\) ）。

命题 3. --- 假设 E 有限维且 A 交换（A = K 或 A = K(i)）。设 \(\Phi\) 是 E 上的埃尔米特型，X 是它关于 E 的基 ( \(x_j\) ) ( \(j = 1, \ldots, n\) ) 的矩阵；对 {[}1, n{]} 的每个子集 H，以 \(X_{\mathrm{H,H}}\) 表示从 X 中删去指标 \(j \notin \mathrm{H}\) 的行与列所得的 X 的子式（第 III 章 § 6 第 3 小节）。

\begin{enumerate}
\def\labelenumi{\alph{enumi})}
\item
  若 \(\Phi\) 非退化且正，则 \(X_{\mathrm{H,H}} > 0\) （对 \([1, n]\) 的每个子集 H）。
\item
  反之，若 \(H_{j} = (1, j)\) 且 \(X_{H_{j}, H_{j}} > 0\) （对 \(j = 1, \ldots, n\) ），则 \(\Phi\) 非退化且正。
\end{enumerate}

先设 \(\Phi\) 非退化且正。诸元素 \((x_{j})\) ( \(j\in\mathbb{H}\) ) 构成 E 的一个子空间 F 的基，而子式 \(X_{\mathrm{H,H}}\) 是限制型 \(\Phi_{\mathrm{F}}\) （即 \(\Phi\) 在 F 上的限制）关于这个基的判别式；今因 \(\Phi_{\mathrm{F}}(x,x)>0\) （对 F 中一切 \(x\neq0\) ），故 \(\Phi_{\mathrm{F}}\) 非退化且正（命题 2 的推论）；因而有 \(X_{\mathrm{H,H}}>0\) （命题 1）。为证明 \(b\) )，注意按 § 6 第 1 小节命题 1 的记号，子式 \(X_{\mathrm{Hj,Hj}}\) 等于 D \(_{j+1,j+1}\) ；因此存在（§ 6 第 1 小节命题 1）E 的正交基 \((e_{j})\) ( \(j=1,\ldots,n\) )，使 \(\Phi(e_{j},e_{j})>0\) （对 \(j=1,\ldots,n\) ）；从而 \(\Phi\) 非退化且正。

注. --- 由标准正交基的存在性可知，两个相同有限维数的向量空间上的两个非退化正埃尔米特型是等价的（§ 1 第 6 小节）。于是设 L 是 A 上一个有限维埃尔米特空间，其度量型非退化且正，V₁ 与 V₂ 是 L 中两个相同维数的线性簇（§ 6 第 6 小节）；由于度量式在 V₁ 与 V₂ 的方向 T₁ 与 T₂ 上的限制、以及在正交子空间 T₁⁰ 与 T₂⁰ 上的限制分别等价，存在 L 的平移空间 T 的酉自同构 \(u\) ，使 \(u(\mathrm{T}_1) = \mathrm{T}_2\) 且 \(u(\mathrm{T}_1^0) = \mathrm{T}_2^0\) ；因而存在 L 的位移 \(v\) 使 \(v(\mathrm{V}_1) = \mathrm{V}_2\) 。设 \((a, b), (a', b')\) 是 L 的点的两个对；要存在 L 的位移 \(v\) 使 \(v(a) = a'\) 且 \(v(b) = b'\) ，必须且只需（用 § 6 第 6 小节的记号）有 \(e(a, b) = e(a', b')\) ；A 的元素 \(\sqrt{e(a, b)}\) （第 VI 章 § 2 第 4 小节）称为埃尔米特空间 L 中 \(a\) 与 \(b\) 的距离。

\subsection{惯性定律。}

定理 1（“惯性定律”）. --- 假设 A 满足本节开头的诸假设，且 E 的维数为有限数 n。设 Φ 是 E 上的一个埃尔米特型。则：

\begin{enumerate}
\def\labelenumi{\alph{enumi})}
\item
  存在 E 的一个分解，把它写成与 E 正交的子空间 \(E^{0}\) 与两个子空间 \(E^{+}\) 与 \(E^{-}\) 的直和，使得 \(\Phi\) 在 \(E^{+}\) （相应地 \(E^{-}\) ）上的限制是正的（相应地，负的）且非退化。
\item
  存在 E 的正交基 \((e_i)_{1\leqslant i\leqslant n}\) ，使得
\end{enumerate}

\[
\Phi (\sum_ {i = 1} ^ {n} \xi_ {i} e _ {i}, \sum_ {i = 1} ^ {n} \eta_ {i} e _ {i}) = \sum_ {i = 1} ^ {s} \xi_ {i} \overline {{\eta}} _ {i} - \sum_ {i = s + 1} ^ {s + t} \xi_ {i} \overline {{\eta}} _ {i}.\tag{2}
\]

\begin{enumerate}
\def\labelenumi{\alph{enumi})}
\setcounter{enumi}{2}
\item
  \(E^{+}\) 的维数 s 与 \(E^{-}\) 的维数 t 对所有满足 a) 中所述条件的直和分解都是相同的；整数 s（相应地，t）是使 \(\Phi\) 在其上的限制为正（相应地，负）且非退化的 E 的子空间 F 的维数的最大值。
\item
  \(\Phi\) 的秩是 \(s + t\) 。
\item
  若 \(\Phi\) 非退化，其指数等于 inf(s, t)（§ 4 第 2 小节定义 2）。
\end{enumerate}

实际上，设 \((x_{i})\) ( \(i = 1, \ldots, n\) ) 是 E 的正交基（§ 6 第 1 小节定理 1）；回顾 \(\Phi(x, x) \in K\) （对一切 \(x \in E\) ）。可以假设 \(\Phi(x_{i}, x_{i}) > 0\) （对 \(i = 1, \ldots, s\) ），\(\Phi(x_{i}, x_{i}) < 0\) （对 \(i = s + 1, \ldots, s + t\) ），\(\Phi(x_{i}, x_{i}) = 0\) （对 \(i = s + t + 1, \ldots, n\) ）。这就证明了 \(a\) )，因为取 \(E^{+}\) （相应地 \(E^{-}\) ）为 \(x_{1}, \ldots, x_{s}\) （相应地 \(x_{s+1}, \ldots, x_{s+t}\) ）生成的子空间，而 \(E^{0}\) 则由 \(x_{s+t+1}, \ldots, x_{n}\) 生成。由此推出 \(b\) )，注意到 \(\Phi(x_{i}, x_{i})\) 具有 \(\rho\bar{\rho}\) 的形式（相应地，

\(-\rho\bar{\rho}) \quad (\rho\in\Lambda^{*}) \quad \text{对} \quad i=1,\ldots,s \quad (\text{resp. } i=s+1,\ldots,s+t).\) 为证明 c)，考察 E 的一个子空间 P，使得 \(\Phi\) 在 P 上的限制是正的且非退化；这时有 \(\mathrm{P}\cap(\mathrm{E}^{-}+\mathrm{E}^{0})=\{0\}\) ，因而和 \(P+E^{-}+E^{0}\) 是直和；我们由此得出 dim \(P\leqslant\dim E^{+}=s\) ，这就证明了 c)。断言 d) 由 a) 立即得出。

最后设 \(\Phi\) 非退化，令 \(q = \inf (s, t)\) ，我们来证明 \(q\) 是 \(\Phi\) 的指数。用 \(b\) ) 的记号，向量 \(e_i + e_{s+i}\) （相应地 \(e_i - e_{s+i}\) ）( \(i = 1, \ldots, q\) ) 生成一个全迷向子空间 F（相应地 F′）。由于 K 的特征为 0，F + F′ 由 \(e_1, \ldots, e_q\) 与 \(e_{s+1}, \ldots, e_{s+q}\) 生成，因而是非迷向的。\(\Phi\) 在子空间 H = (F + F′)⁰ 上的限制于是是正（或负）的且非退化，故 H 不含 \(\neq 0\) 的迷向向量。由于 F⁰ 包含 F 与 H，且 dim(F + H) = codim F′ = codim F，有 F⁰ = F + H。于是与 F 正交的迷向向量 z 必属于 F，因为在直和 F + H 中，z 在 H 中的分量是迷向的。因此 F 是一个极大全迷向子空间，这就证明了 e)。

定义 2. --- 用定理 1 的记号，称偶 \((s,t)\) 为 \(\Phi\) 的符号差。

\subsection{型关于正埃尔米特型的约化。}

在本节中，我们将假设 E 在 A 上的维数为有限数 n，并以 Φ 表示 E 上的一个非退化正埃尔米特型。

由于从 E 到 \(\mathbf{E}^*\) 中的与 \(\Phi\) 相伴的线性映射都是双射，故对 E 中任意元素 \(x, y, z\) （ \(y \neq 0\) ），存在 E 的恰好一个元素 \(t\) 使 \(\Phi(x, z) = \overline{\Phi(t, y)}\) 。特别地，若 \(\mathrm{A} = \mathrm{K}\) 或 \(\mathrm{A} = \mathrm{K}(i)\) ，且 \(u\) 是 E 到自身的（关于 J 的）半线性映射（第 II 章，附录 I，第 1 小节），则对每个 \(x \in \mathrm{E}\) ，存在恰好一个元素 \(u^*(x)\) ，使得对一切 \(y \in \mathrm{E}\) 有 \(\Phi(x, u(y)) = \overline{\Phi(u^*(x), y)}\) 。立即可以看出，\(u^*\) 是 E 到自身的（关于 J 的）半线性映射；称之为 \(u\) 的伴随。

注记. --- 1) 当 J 是恒等映射时，我们重新得到 § 1 第 8 小节中定义的同态的伴随概念。

\begin{enumerate}
\def\labelenumi{\arabic{enumi})}
\setcounter{enumi}{1}
\tightlist
\item
  仍假设 A = K 或 A = K(i)。设 E \(^{j}\) 是 § 1 第 2 小节定义 5 中定义的右 A-模。型 Φ 是 E × E \(^{j}\) 上的双线性型，Φ \(^{j}\) 是 E \(^{j}\) × E 上的双线性型，而 u 是从 E 到 E \(^{j}\) 中的 A-线性映射。这个同态在 § 1 第 8 小节意义下的伴随是从 E 到 E \(^{j}\) 中的一个 A-线性映射；我们立即看出它与上面定义的映射 u* 重合。
\end{enumerate}

我们称 E 的自同态 \(u\) 是（关于 \(\Phi\) 的）正规自同态，如果有 \(uu^{*}=u^{*}u\) 。
